\documentclass[10pt,a4paper]{amsart}
\usepackage{amsmath,amssymb,amsthm,graphicx}
% Retained preamble for the Stacks Project chapter {CH} (shared project preamble).
% The chapter file pinned-original-source.tex only carries \input{{preamble}} on its first line,
% so no preamble text lives inside the chapter itself.  The Stacks Project uses one single shared
% preamble for every chapter: the snapshot stages it three times (sources/stacks-chow/preamble.tex,
% sources/stacks-etale-cohomology/preamble.tex and focused-seeds/AGP/sources/stacks-snapshot-20261011/
% preamble.tex) and those copies are identical in size.  The block below is copied mechanically,
% line for line, from sources/stacks-chow/preamble.tex lines 14-18 and 154-217: the Xy-pic setup and
% the native theorem declarations together with the mathematical macros.  Nothing was added, removed
% or redefined; the wrapper class and the amsmath/amssymb/amsthm/graphicx base packages are supplied
% by the assembler.
% For commutative diagrams we use Xy-pic
% Chapter: stacks-cohomology.tex -- Cohomology of Sheaves.
\usepackage[all]{xy}

% We use 2cell for 2-commutative diagrams.
\xyoption{2cell}

% Theorem environments.
%
\theoremstyle{plain}
\newtheorem{theorem}[subsection]{Theorem}
\newtheorem{proposition}[subsection]{Proposition}
\newtheorem{lemma}[subsection]{Lemma}

\theoremstyle{definition}
\newtheorem{definition}[subsection]{Definition}
\newtheorem{example}[subsection]{Example}
\newtheorem{exercise}[subsection]{Exercise}
\newtheorem{situation}[subsection]{Situation}

\theoremstyle{remark}
\newtheorem{remark}[subsection]{Remark}
\newtheorem{remarks}[subsection]{Remarks}

\numberwithin{equation}{subsection}

% Macros
%
\def\lim{\mathop{\mathrm{lim}}\nolimits}
\def\colim{\mathop{\mathrm{colim}}\nolimits}
\def\Spec{\mathop{\mathrm{Spec}}}
\def\Hom{\mathop{\mathrm{Hom}}\nolimits}
\def\Ext{\mathop{\mathrm{Ext}}\nolimits}
\def\SheafHom{\mathop{\mathcal{H}\!\mathit{om}}\nolimits}
\def\SheafExt{\mathop{\mathcal{E}\!\mathit{xt}}\nolimits}
\def\Sch{\mathit{Sch}}
\def\Mor{\mathop{\mathrm{Mor}}\nolimits}
\def\Ob{\mathop{\mathrm{Ob}}\nolimits}
\def\Sh{\mathop{\mathit{Sh}}\nolimits}
\def\NL{\mathop{N\!L}\nolimits}
\def\CH{\mathop{\mathrm{CH}}\nolimits}
\def\proetale{{pro\text{-}\acute{e}tale}}
\def\etale{{\acute{e}tale}}
\def\QCoh{\mathit{QCoh}}
\def\Ker{\mathop{\mathrm{Ker}}}
\def\Im{\mathop{\mathrm{Im}}}
\def\Coker{\mathop{\mathrm{Coker}}}
\def\Coim{\mathop{\mathrm{Coim}}}

% Boxtimes
%
\DeclareMathSymbol{\boxtimes}{\mathbin}{AMSa}{"02}

%
% Macros for moduli stacks/spaces
%
\def\QCohstack{\mathcal{QC}\!\mathit{oh}}
\def\Cohstack{\mathcal{C}\!\mathit{oh}}
\def\Spacesstack{\mathcal{S}\!\mathit{paces}}
\def\Quotfunctor{\mathrm{Quot}}
\def\Hilbfunctor{\mathrm{Hilb}}
\def\Curvesstack{\mathcal{C}\!\mathit{urves}}
\def\Polarizedstack{\mathcal{P}\!\mathit{olarized}}
\def\Complexesstack{\mathcal{C}\!\mathit{omplexes}}
% \Pic is the operator that assigns to X its picard group, usage \Pic(X)
% \Picardstack_{X/B} denotes the Picard stack of X over B
% \Picardfunctor_{X/B} denotes the Picard functor of X over B
\def\Pic{\mathop{\mathrm{Pic}}\nolimits}
\def\Picardstack{\mathcal{P}\!\mathit{ic}}
\def\Picardfunctor{\mathrm{Pic}}
\def\Deformationcategory{\mathcal{D}\!\mathit{ef}}

% --- appended 2026-10-12 (Chief): the project own definitions of the reference/slogan/history
% display environments, copied mechanically from the shared Stacks preamble lines 8-12
% (\usepackage{verbatim} plus the three \newenvironment lines, which the chapter preamble keeps
% verbatim from the project preamble). Nothing added, removed or redefined: these environments
% carry no mathematical content and the project itself typesets them as comments. Without these
% three lines a counted statement whose own text uses them fails with Environment slogan undefined.
\usepackage{verbatim}
\newenvironment{reference}{\comment}{\endcomment}
\newenvironment{slogan}{\comment}{\endcomment}
\newenvironment{history}{\comment}{\endcomment}
\title{Filtration of a finitely generated subsheaf of the constant sheaf Z by extensions of extensions by zero}
\author{The Stacks Project Authors}
\date{stacks-project:cohomology.tex:local-snapshot-20261011; GNU FDL 1.2 or later}
\begin{document}
\maketitle
\noindent\textit{Editorial note.} All delivered mathematical text is the author's own, copied from the pinned Stacks Project chapter file cohomology.tex; no statement or proof step was rewritten, reordered or removed.\par\medskip
\noindent\textit{Editorial note.} This card was re-admitted after a gate audit: it had been excluded only because the chapter's retained preamble lacks the reference/slogan/history environments, and the author's own preamble defines them as comment environments. The card therefore carries a per-card preamble copy in its task directory, consisting of the chapter retained preamble copied verbatim plus the project's own three display environments copied verbatim from sources/stacks-chow/preamble.tex lines 7-12. Nothing was added or redefined by the assembler side.\par\medskip
\noindent\textit{Editorial note.} The author's own bibliographic remark, which sits in a reference environment at the head of the statement and cites the original paper, is preserved in the source span and is not typeset because the author's own preamble comments reference environments out; no bibliography excerpt is therefore needed.\par\medskip
\noindent\textit{Editorial note.} The counted spans carry no navigation stub other than those already listed and no bibliography citation reaches the compiled document.\par\medskip
\noindent\textit{Editorial note.} The label lemma-subsheaf-of-constant-sheaf is defined only in cohomology.tex, so no cross-chapter restatement had to be weighed for this card.\par\medskip
\noindent\textit{Allowed standard prerequisites (editorial).}\begin{itemize}
\item Subsheaf generated by local sections. Modules, Definition modules-definition-generated-by-local-sections, and the description of the subsheaf generated by finitely many sections, used throughout proof lines 3504-3526. Supplies the sheaves whose filtration is constructed.
\item Extension by zero along the inclusion of an open subspace. Sheaves, the extension by zero functor named in the statement at lines 3496-3500. Supplies the terms of the short exact sequences.
\item Sections of the constant sheaf Z are locally constant integer valued functions. Standard, used at proof lines 3505-3508. Supplies the decomposition that makes the generators constant.
\end{itemize}\medskip
\begin{lemma}
\label{lemma-subsheaf-of-constant-sheaf}
\begin{reference}
\cite[Page 168]{Tohoku}.
\end{reference}
Let $X$ be a topological space such that the intersection of any
two quasi-compact opens is quasi-compact. Let
$\mathcal{F} \subset \underline{\mathbf{Z}}$
be a subsheaf generated by finitely many sections over quasi-compact opens.
Then there exists a finite filtration
$$
(0) = \mathcal{F}_0 \subset \mathcal{F}_1 \subset \ldots \subset
\mathcal{F}_n = \mathcal{F}
$$
by abelian subsheaves such that for each $0 < i \leq n$
there exists a short exact sequence
$$
0 \to j'_!\underline{\mathbf{Z}}_V \to j_!\underline{\mathbf{Z}}_U \to
\mathcal{F}_i/\mathcal{F}_{i - 1} \to 0
$$
with $j : U \to X$ and $j' : V \to X$ the inclusion of quasi-compact opens
into $X$.
\end{lemma}

\begin{proof}
Say $\mathcal{F}$ is generated by the sections $s_1, \ldots, s_t$ over the
quasi-compact opens $U_1, \ldots, U_t$. Since $U_i$ is quasi-compact and
$s_i$ a locally constant function to $\mathbf{Z}$ we may assume, after
possibly replacing $U_i$ by the parts of a finite decomposition into open
and closed subsets, that $s_i$ is a constant section.
Say $s_i = n_i$ with $n_i \in \mathbf{Z}$. Of course we can remove
$(U_i, n_i)$ from the list if $n_i = 0$. Flipping signs if necessary
we may also assume $n_i > 0$. Next, for any subset $I \subset \{1, \ldots, t\}$
we may add $\bigcap_{i \in I} U_i$ and $\gcd(n_i, i \in I)$ to the list.
After doing this we see that our list $(U_1, n_1), \ldots, (U_t, n_t)$
satisfies the following property:
For $x \in X$ set $I_x = \{i \in \{1, \ldots, t\} \mid x \in U_i\}$.
Then $\gcd(n_i, i \in I_x)$ is attained by $n_i$ for some $i \in I_x$.

\medskip\noindent
As our filtration we take $\mathcal{F}_0 = (0)$ and
$\mathcal{F}_n$ generated by the sections $n_i$ over $U_i$ for those
$i$ such that $n_i \leq n$. It is clear that
$\mathcal{F}_n = \mathcal{F}$ for $n \gg 0$. Moreover, the quotient
$\mathcal{F}_n/\mathcal{F}_{n - 1}$ is generated by the section
$n$ over $U = \bigcup_{n_i \leq n} U_i$ and the kernel of the map
$j_!\underline{\mathbf{Z}}_U \to \mathcal{F}_n/\mathcal{F}_{n - 1}$
is generated by the section $n$ over $V = \bigcup_{n_i \leq n - 1} U_i$.
Thus a short exact sequence as in the statement of the lemma.
\end{proof}

\end{document}
